% Generated by Listas/make_statements.py from pratica-lecture-04.md, item 10. Do not edit.
\textbf{Source.} Tadelis (2013), \emph{Game Theory: An Introduction}, Exercise 8.7 ``Debt and Repayment'', p.~171 {[}PDF 190{]}. Reproduced verbatim. Answer in the source: none in the book (the preface, p.~xiii, says about half are solved on Princeton UP's website; not checked).



\begin{quote}

\textbf{Debt and Repayment:} A project costing \$100 yields a gross return of \$110. A lender (player 1) is approached by a debtor (player 2) requesting a standard loan contract to complete the project. If the lender chooses \emph{not to offer} a loan, then both parties earn nothing. If the lender chooses to \emph{offer} a loan of \$100, the debtor can realize the project's gains and is obliged by contract to repay \$105. For simplicity assume that money is continuous and that the debtor can choose to return any amount of money \(x \leq 110\). Ignore the time value of money. Assume first that no legal system is in place that can compel the debtor to repay, so that default on the loan (less than full repayment) carries no repercussions for the debtor.



\begin{enumerate}

\def\labelenumi{\alph{enumi}.}

\item

  Model this as an extensive-form game tree as best you can and find a subgame-perfect equilibrium of this game. Is it unique?

\item

  Now assume that there is a legal system in place that allows the lender to \emph{voluntarily} choose whether or not to sue when the debtor defaults and repays an amount \(x < 105\). Furthermore assume that it is \emph{costless} to use the legal system (it is supplied by the state), and if the lender sues a debtor that defaulted, the lender will get the \$105 repaid in full. After paying the lender, the borrower will pay a fine of \$5 to the court, above and beyond the repayment. Model this as an extensive-form game tree as best you can and find a subgame-perfect equilibrium of this game. Is it unique?

\item

  Are there Nash equilibria in the game described in (b) that are not subgame-perfect equilibria?

\item

  Now assume that using the legal system is \emph{costly:} if the lender sues, he pays lawyers a \emph{legal fee} of \$105 (this is the lawyers' price; it is unrelated to the contract). The rest of the game proceeds as before. (If the lender sues a debtor that defaulted, the lender will get repaid in full; after paying the lender, the borrower will pay a fine of \$5 above and beyond the repayment.) Model this as an extensive-form game tree as best you can and find a subgame-perfect equilibrium of this game. Is it unique?

\item

  Are there Nash equilibria in the game described in (d) that are not subgame-perfect equilibria?

\item

  Now assume that a change in the law is proposed: upon default, if a debtor is sued he has to first repay the lender \$105 and then pay the legal fees of \$105 above and beyond repayment of the loan, and no extra fine is imposed. Should the lender be willing to pay to have this change in the law enacted? If so, how much?

\item

  If you were a social planner, would you implement the suggested law?

\end{enumerate}

\end{quote}

